%=========================
% Preamble/style.tex
% Style chung: màu sắc, header/footer, caption, box, hyperlink, clever reference
%=========================

%========== Tên thành phần tiếng Việt ==========
\renewcommand{\contentsname}{Mục lục}
\renewcommand{\listfigurename}{Danh mục hình}
\renewcommand{\listtablename}{Danh mục bảng}
\renewcommand{\figurename}{Hình}
\renewcommand{\tablename}{Bảng}
\renewcommand{\refname}{Tài liệu tham khảo}
\floatname{algorithm}{Thuật toán}
\renewcommand{\algorithmicrequire}{\textbf{Đầu vào:}}
\renewcommand{\algorithmicensure}{\textbf{Đầu ra:}}

%========== Typography & paragraph ==========
\onehalfspacing
\setlength{\parindent}{1.2cm}
\setlength{\parskip}{3pt}
\setlist{itemsep=3pt, topsep=4pt}
\renewcommand{\arraystretch}{1.18}

%========== Header / footer ==========
\pagestyle{fancy}
\fancyhf{}
\setlength{\headheight}{40pt}
\fancyhead[L]{%
  \begingroup
  \setstretch{1}
  \normalsize
  \renewcommand{\arraystretch}{1}
  \begin{tabular}{@{}rl@{}}
    \begin{picture}(25,15)(0,0)
      \put(0,-8){%
        \includegraphics[width=10mm,height=10mm]{Images/\LogoFile}%
      }
    \end{picture}
    &
    \begin{tabular}{@{}l@{}}
      \textbf{Đại học Quốc gia TP. Hồ Chí Minh}\\
      \textbf{Trường Đại học Bách Khoa - Khoa KH\&KT Máy tính}
    \end{tabular}
  \end{tabular}
  \endgroup
}
\fancyfoot[L]{\scriptsize Báo cáo Đồ án chuyên ngành}
\fancyfoot[R]{%
  \if\relax\thepage\relax
  \else
    \scriptsize Trang~\thepage/\pageref*{LastPage}%
  \fi
}
\renewcommand{\headrulewidth}{0.3pt}
\renewcommand{\footrulewidth}{0.3pt}

%========== Tiêu đề kiểu chương trong article ==========
% Section vẫn đánh số 1, 2, 3 để subsection là 1.1, 1.2...
% Nhưng heading hiển thị là "Chương 1 -- Tên chương".
\titleformat{\section}
  {\clearpage\Large\bfseries\color{HCMUTDarkBlue}}
  {Chương~\thesection}{0.8em}{--\ }
\titleformat{\subsection}
  {\large\bfseries\color{HCMUTBlue}}
  {\thesubsection}{0.8em}{}
\titleformat{\subsubsection}
  {\normalsize\bfseries\color{HCMUTDarkBlue}}
  {\thesubsubsection}{0.8em}{}
\titlespacing*{\section}{0pt}{0pt}{10pt}
\titlespacing*{\subsection}{0pt}{12pt}{6pt}
\titlespacing*{\subsubsection}{0pt}{10pt}{4pt}

%========== Mục lục ==========
\renewcommand{\cftsecleader}{\cftdotfill{\cftdotsep}}
\setlength{\cftbeforesecskip}{4pt}

%========== Caption ==========
\captionsetup{
  labelfont={bf,color=HCMUTDarkBlue},
  textfont=it,
  font=small,
  justification=centering
}
\captionsetup[table]{position=top}

%========== Link ==========
\usepackage[unicode]{hyperref}
\hypersetup{
  colorlinks=true,
  linkcolor=HCMUTDarkBlue,
  citecolor=HCMUTDarkBlue,
  urlcolor=HCMUTBlue,
  filecolor=HCMUTBlue
}
\usepackage{bookmark}
\usepackage[nameinlink,noabbrev]{cleveref}

% Tên tham chiếu tiếng Việt cho cleveref.
\crefname{section}{chương}{chương}
\Crefname{section}{Chương}{Chương}
\crefname{subsection}{mục}{mục}
\Crefname{subsection}{Mục}{Mục}
\crefname{subsubsection}{mục}{mục}
\Crefname{subsubsection}{Mục}{Mục}
\crefname{figure}{hình}{hình}
\Crefname{figure}{Hình}{Hình}
\crefname{table}{bảng}{bảng}
\Crefname{table}{Bảng}{Bảng}
\crefname{algorithm}{thuật toán}{thuật toán}
\Crefname{algorithm}{Thuật toán}{Thuật toán}
\crefname{lstlisting}{mã nguồn}{mã nguồn}
\Crefname{lstlisting}{Mã nguồn}{Mã nguồn}

%========== tcolorbox: box dùng trong báo cáo ==========
\tcbset{
  enhanced,
  boxrule=0.6pt,
  arc=2mm,
  left=2mm,
  right=2mm,
  top=1.5mm,
  bottom=1.5mm,
  fonttitle=\bfseries,
  breakable
}

\newtcolorbox{infobox}[1]{
  title={#1},
  colback=HCMUTLightBlue,
  colframe=HCMUTBlue
}

\newtcolorbox{notebox}[1]{
  title={#1},
  colback=HCMUTLightBlue,
  colframe=HCMUTDarkBlue,
  coltitle=white,
  colbacktitle=HCMUTDarkBlue
}

\newtcolorbox{warningbox}[1]{
  title={#1},
  colback=WarningYellow,
  colframe=orange!70!black
}

\newtcolorbox{requirementbox}[1]{
  title={#1},
  colback=SuccessGreen,
  colframe=green!45!black
}

\newtcolorbox{riskbox}[1]{
  title={#1},
  colback=DangerRed,
  colframe=red!60!black
}

%========== Macro tiện ích ==========
\newcommand{\filepath}[1]{\texttt{\detokenize{#1}}}
\newcommand{\inlinecode}[1]{\texttt{\detokenize{#1}}}
\newcommand{\term}[1]{\textbf{#1}}
\newcommand{\todoitem}[1]{\textcolor{red!70!black}{\textbf{TODO:} #1}}
\newcommand{\sourceplaceholder}{\textcolor{red!70!black}{[Cần trích dẫn nguồn]}}


% Tạo heading không đánh số, không đưa vào mục lục, không hiện số trang.
\newcommand{\frontmatterchapter}[1]{%
  \clearpage
  \section*{#1}
  \thispagestyle{empty}
}

%========== Cấu hình siunitx ==========
\sisetup{
  detect-all,
  group-separator={,},
  group-minimum-digits=4
}

%feedback 
\newcommand{\commentlines}[1]{%
  \par\noindent
  \begingroup
  \color{black!35}%
  \count0=0
  \loop\ifnum\count0<#1
    \noindent\rule{\textwidth}{0.4pt}\par\vspace{0.75cm}%
    \advance\count0 by 1
  \repeat
  \endgroup
}